\chapter{Sachlich und modal}

\section{Differenzen des Gegenstandes, Differenzen des Anerkennens}

Die drei Asymmetrien, Grenze, Richtung und Maß, haben einen gemeinsamen Grund, und Brentano hat ihn benannt. Räumliche Differenzen sind Differenzen des Gegenstandes. Was hier ist, ist ein anderes als was dort ist, so wie Rot ein anderes ist als Blau. Zeitliche Differenzen sind keine Differenzen des Gegenstandes. Dass etwas vergangen, gegenwärtig oder künftig ist, ist keine Bestimmung an ihm, sondern eine Weise, wie es vorgestellt und anerkannt wird: ein Modus. Vergangen und künftig sind Modi in obliquo, gegenwärtig ist der Modus in recto. Und diese Modi sind komparative Relativa zum Sprechenden: Was für mich vergangen ist, ist es im Verhältnis zu meinem Jetzt.\footnote{Brentano, Untersuchungen, Kapitel VII, Sachliche und modale Differenzen, und VIII, Das Zeitliche als Relatives.}

Brentano prüft, was geschähe, wenn es anders wäre. Gäbe es nur einen einzigen Präteritalmodus, so wäre das Vergangene vom Gegenwärtigen unterschieden wie ein Ort vom anderen; die Zeit wäre dann \zitat{ein Topoid von einer Dimension}, ein eindimensionaler Raum, und das Nacheinander wäre ein Nebeneinander.\footnote{Brentano, Untersuchungen, Kapitel IV, zum Präteritalmodus.} Dass die Zeit kein Topoid ist, liegt daran, dass die Modi kontinuierlich variieren, und dass sie eine einzige Dimension hat, ist keine Zufälligkeit wie beim Raum: Ein Chronoid von anderer Dimensionszahl ist \zitat{unmittelbar absurd}, während Topoide von beliebiger Dimensionszahl denkbar sind. Der Raum könnte anders sein, als er ist; die Zeit nicht.

Die Überordnung, die daraus folgt, spricht Brentano aus: Der Begriff des Zeitlichen \zitat{verhält sich zu dem des Räumlichen als ein übergeordneter} und begreift Spezies unter sich, die aller räumlichen Kontinuität ermangeln.\footnote{Brentano, Untersuchungen, S.\,213.} Alles Reale ist zeitlich; nicht alles Reale ist räumlich. Und noch einmal die Kreuzung, jetzt an Hier und Jetzt: \zitat{Der Ortspunkt ist Grenze für ein Räumliches in vielen, ja vielleicht in allen denkbaren räumlichen Richtungen, aber niemals Grenze in einer der beiden zeitlichen Richtungen, und vom Zeitpunkt gilt das Umgekehrte.}\footnote{Brentano, Untersuchungen, S.\,208.}

Brentano hat diese Lehre später abgeschwächt und auch in der Zeit relative sachliche Differenzen zugelassen, den Abstand als Bestimmung des Gegenstandes. Aber der Kern bleibt: Der Abstand in der Zeit ist ein Abstand zwischen Weisen, in denen derselbe Gegenstand anerkannt wird. Man kann sich denselben Ton, dieselbe Farbe, denselben Zustand nach beliebig langer Zeit wiederholt denken, ohne dass sich etwas an ihm geändert hätte. Man kann sich nicht denselben Punkt an zwei Orten denken.

\section{Warten und Gehen}

Chiba gibt der modalen Asymmetrie eine Form, die ohne Brentanos Terminologie auskommt: \zitat{Um die Zukunft zu erfahren, muss man nur warten. Dergleichen gilt nicht für den Raum.}\footnote{Chiba, Kants Ontologie der raumzeitlichen Wirklichkeit, Kapitel zur raumzeitlichen Wirklichkeit.} Um das Dort zu erfahren, muss man hingehen, und Gehen ist eine Tätigkeit, die im Raum eine Richtung wählt. Warten ist keine Tätigkeit, die eine Richtung wählt. Die Zukunft kommt von selbst; das Dort kommt nicht von selbst. Das ist die Einsinnigkeit, die Reichenbach nicht definieren konnte, in der Sprache der Erfahrung. Sie ist kein Merkmal der Ereignisse, sondern der Weise, wie sie sich einstellen.

Chiba fasst die raumzeitlichen Gegenstände als solche, die räumlich \emph{und/oder} zeitlich sind, und die Zeit ist die weitere Klasse: Es gibt Zeitliches ohne Räumliches, ein Gedanke etwa, aber kein Räumliches ohne Zeitliches. Das ist Brentanos Überordnung bei Kant. Und Chiba gibt dem Realismus, den er bestreitet, einen Namen, der die Sache trifft: Er sei ein \zitat{zeit-transzendenter Standpunkt}, ein Blick von nirgendwann. Die Existenz eines Gegenstandes ist für ihn dessen Erkennbarkeit; leere Zeit und leerer Raum sind nichts Wirkliches, weil sie nichts sind, was erkannt werden könnte.\footnote{Chiba, Kants Ontologie, Kapitel zum Anti-Realismus und zum Existenzbegriff.}

Der Ausdruck \zitat{zeit-transzendent} ist genauer, als er zunächst scheint. Der Realismus, der die Welt als das nimmt, was sie unabhängig von jedem Standpunkt ist, muss sich einen Blick denken, der nicht jetzt ist. Einen Blick, der nicht hier ist, kann man sich leicht denken: Man geht einfach woanders hin, oder man denkt sich alle Orte zugleich. Einen Blick, der nicht jetzt ist, kann man sich nur denken, indem man die Zeit zum Raum macht, in dem alle Jetzt nebeneinander liegen. Der zeit-transzendente Standpunkt ist also der Standpunkt, von dem aus die Zeit als Raum erscheint.

\section{Die Physik und das Nichts}

Koch hat dieselbe Sache von der Physik her gesagt. Die Physik abstrahiert von der Subjektivität. Sie hat die Kontinuität der Zeit mittels analytischer Geometrie in diskrete reelle Zahlen aufgelöst; die Zeit wird zum Parameter $t$, der wie eine Koordinate behandelt wird. Damit hat sie genau das von der Zeit behalten, worin die Zeit dem Raum gleicht, die Ordnung, und das weggelassen, worin sie sich unterscheidet, den Modus. Koch nennt die Subjektivität das Nichts, den Ursprung der Negation, des Werdens, der Bewegung, und die Physik den Versuch, die Welt unter Absehung vom Nichtigen zu begreifen, ein eleatisches Unternehmen.\footnote{Koch, Erste Philosophie, Abschnitt zur Physik und zur Subjektivität.}

Das ist keine Kritik an der Physik. Es ist die Feststellung, was sie tut. Reichenbach selbst hat es festgestellt, als er die Wahrnehmungsordnung von der Ereignisordnung unterschied: Die Wahrnehmungen bilden eine eindimensionale Kette, die objektiven Ereignisse eine Mannigfaltigkeit mit einer Dimension mehr, in der nichts \zitat{jetzt} ist.\footnote{Reichenbach, Raum-Zeit-Lehre, \S\,45.} Und er hat die Ich-Zeit, in der das \zitat{ewige Jetzt} vorkommt, ausdrücklich vertagt. Sie ist im Buch nicht wiedergekommen.

Was hier zusammenkommt, ist die philosophische Form der physikalischen Asymmetrie. Hegel hat die Zeit das angeschaute Werden genannt, das Umschlagen von Sein in Nichts und Nichts in Sein; Bonsiepen hat gezeigt, dass Hegel den Raum dagegen substanzmetaphysisch, als ruhendes Sein denkt.\footnote{Bonsiepen, Hegels Raum-Zeit-Lehre, zu den Nachschriften und zur Genese.} Die Zeit gehört auf die Seite des Nichts, des Modus, der Subjektivität; der Raum auf die Seite des Seins, des Gegenstandes, der Sache. Die Physik, die vom Nichts absieht, muss die Zeit deshalb nach dem Raum behandeln, und sie tut es mit Recht, denn sie will die Sache. Was sie dabei verliert, ist nicht ein Teil der Sache, sondern der Modus, in dem die Sache je gegenwärtig ist. Und dieser Modus ist, was die Zeit von allem unterscheidet, das man nebeneinanderlegen kann.
